\appendix

\setcounter{table}{0}
\renewcommand{\thetable}{A\arabic{table}}
\renewcommand{\theHtable}{appendix.\arabic{table}}
\setcounter{figure}{0}
\renewcommand{\thefigure}{A\arabic{figure}}
\renewcommand{\theHfigure}{appendix.\arabic{figure}}
\setcounter{algorithm}{0}
\renewcommand{\thealgorithm}{A\arabic{algorithm}}
\renewcommand{\theHalgorithm}{appendix.\arabic{algorithm}}

\section{附录}
\label{app:appendix}

\subsection{数据卡}
\label{app:datacard}

\begin{table}[ht]
\centering
\small
\caption{TwinRouterBench 静态赛道的数据卡。它遵循为 LLM 基准发布调整后的 \citet{bender2018data} 和 \citet{gebru2021datasheets} 数据声明惯例。}
\label{tab:datacard}
\begin{tabularx}{\linewidth}{p{0.28\linewidth} X}
\toprule
字段 & 内容 \\
\midrule
\textbf{数据集名称} & TwinRouterBench v1.0（静态赛道） \\
\textbf{任务类型} & 步级 LLM 路由（在层级 ID 0--3 上进行的四分类） \\
\textbf{规模} & 来自 520 个轨迹实例的 970 条记录 \\
\textbf{源工作负载} & SWE-bench（Apache-2.0）、BFCL（CC BY 4.0）、mtRAG（IBM Research；见原始许可证）、QMSum（MIT）、PinchBench~\citep{pinchbench}（MIT；其 53 个任务中 12 个任务派生的 48 条记录被纳入本发布版本，并受下述包级 Apache-2.0 许可证覆盖） \\
\textbf{来源} & 从强模型成功执行下的智能体轨迹中提取消息前缀；层级标签由贪心顺序锁定降级搜索和执行验证产生 \\
\textbf{标签语义} & 逐步最低成本充分层级的估计：在 11 模型池、由 Opus 初始化的顺序锁定降级搜索，以及 \S\ref{sec:dataset-construction} 的级联协议下，通过固定步数的端到端轨迹通过进行验证；即模型池中能够正确处理当前步骤的最低层级，并通过人工审计交叉检查 \\
\textbf{许可证} & Apache-2.0 \\
\textbf{预期用途} & 训练和评测步级 LLM 路由器；对智能体部署成本降低进行基准评测 \\
\textbf{范围外用途} & 宣称绝对路由最优性；在未重新标注的情况下部署至与发布层级映射显著不同的模型池；作为通用 SWE-bench 或 BFCL 能力基准使用 \\
\textbf{潜在偏差} & \texttt{high} 层级中过度代表智能体代码修复（SWE-bench）；BFCL/mtRAG/QMSum 在 \texttt{low} 层级近乎饱和；仅有英文内容；标签依赖模型池和 harness \\
\textbf{维护计划} & 新的层级池清单和评分协议版本将通过 GitHub releases 跟踪；层级映射和定价均版本化；当前发布为 v1.0 \\
\textbf{人工监督} & 对 BFCL、SWE-bench 和 PinchBench 的随机步骤样本进行最终人工审计（总计 64 步，约 $\sim$10\%）；一个 SWE-bench 步骤被标记为可继续降级，且其验证层级在发布前下调一级；审计样本中未发现其他仍可降级的标签 \\
\textbf{发布 URL} & 公开代码与数据包：\url{https://github.com/CommonstackAI/TwinRouterBench} \\
\bottomrule
\end{tabularx}
\end{table}

许可证元数据遵循发布日时的上游文档；用户在重新分发派生制品前应核实上游条款。

\subsection{语料库概览}
\label{app:corpus-overview}

\begin{table}[ht]
\centering
\small
\begin{adjustbox}{max width=\linewidth}
\begin{tabular}{lrrrrl}
\toprule
基准 & 实例数 & 步骤数 & 每条轨迹步骤数 & 前缀 token & 场景 \\
 & & & （中位数） & （中位数） & \\
\midrule
SWE-bench      & 40  & 336 & 9.0 & $\sim$5{,}300 & GitHub issue 代码修复（多步智能体） \\
BFCL           & 130 & 248 & 1.0 & $\sim$1{,}600 & 函数调用（单轮与多轮） \\
mtRAG          & 193 & 193 & 1   & $\sim$1{,}900 & 多轮检索增强问答 \\
QMSum          & 145 & 145 & 1   & $\sim$3{,}000 & 面向查询的会议摘要 \\
PinchBench     & 12  & 48  & 4   & $\sim$10{,}500 & 含 23 个任务的通用智能体套件 \\
\midrule
\textbf{总计} & \textbf{520} & \textbf{970} & & & \\
\bottomrule
\end{tabular}
\end{adjustbox}
\caption{语料库概览。前缀 token 的中位数使用 Anthropic tokenizer 在路由器可见的 \texttt{messages} 字段上计算。}
\label{tab:corpus-overview}
\end{table}

\subsection{目标层级分布与顺序锁定搜索}
\label{app:tier-alg}

为节省篇幅，我们将工作负载目标层级计数（正文 \S\ref{sec:benchmark-overview} 引用）与顺序锁定降级搜索伪代码（正文 \S\ref{sec:dataset-construction} 引用）置于本附录。

\begin{table}[ht]
\centering
\small
\begin{tabular}{lrrrr}
\toprule
基准 & \texttt{low} & \texttt{mid} & \texttt{mid\_high} & \texttt{high} \\
\midrule
BFCL       & 239 & 8  & 1  & 0 \\
mtRAG      & 183 & 8  & 1  & 1 \\
QMSum      & 132 & 10 & 3  & 0 \\
PinchBench & 41  & 3  & 3  & 1 \\
SWE-bench  & 94  & 33 & 41 & 168 \\
\midrule
总计      & 689 & 62 & 49 & 170 \\
\bottomrule
\end{tabular}
\caption{按工作负载划分的目标层级分布。}
\label{tab:tier-distribution}
\end{table}

\begin{algorithm}[ht]
\caption{顺序锁定降级搜索（每条轨迹）}
\label{alg:seq-lock}
\scriptsize
\begin{algorithmic}[1]
\REQUIRE 轨迹 $\tau$、提示 $h_1,\ldots,h_N$、harness $\mathcal{H}$、层级 $T_0<T_1<T_2<T_3$
\STATE 对所有 $i$，令 $\text{locked}[i]\leftarrow T_3$
\FOR{$i=1$ \textbf{到} $N$}
  \IF{$h_i=$ 不可降级}
    \STATE \textbf{继续}
  \ENDIF
  \FOR{$t$ 从 $h_i$ \textbf{到} $T_2$}
    \STATE 运行 $\mathcal{H}$：步骤 $<i$ 已锁定，步骤 $i$ 位于层级 $t$，步骤 $>i$ 位于层级 $T_3$
    \IF{轨迹通过}
      \STATE $\text{locked}[i]\leftarrow t$；\textbf{终止循环}
    \ENDIF
  \ENDFOR
\ENDFOR
\RETURN $\text{locked}$
\end{algorithmic}
\end{algorithm}

\subsection{制品结构}
\label{app:artifact}

审稿包包含检查基准和重新计算报告表格所需的静态 JSONL、锁定清单、评分代码、动态切分与紧凑动态摘要。预期布局如下：

\begin{center}
\begin{minipage}{0.95\linewidth}
\scriptsize
\begin{verbatim}
data/static/question_bank.jsonl
data/static/manifest.json
data/dynamic/model_pool.json
data/dynamic/model_pricing.json
data/dynamic/tier_to_model.json
data/dynamic/ttl_policy.json
main/eval/
main/metrics.py
README.md
\end{verbatim}
\end{minipage}
\end{center}

静态文件包含 970 条记录，单条记录只公开层级标签，不公开提供商模型 ID。动态评分使用 \texttt{data/dynamic/} 下的锁定模型池、定价、层级映射和 TTL 策略文件；留出的 SWE-bench ID 和聚合输出随发布制品一并提供。README 提供静态评分器和动态评分 CLI 的安装与 smoke-test 命令。对于 E\&D 赛道投稿，发布包还包括 \texttt{metadata/croissant.json}，其中含 Croissant 核心字段与 Responsible AI 字段、许可证和源工作负载元数据，以及静态评分器和动态摘要评分器的可执行 smoke-test 脚本。

\subsection{静态模型池}
\label{app:model-pool}

静态赛道使用一个固定的、由 11 个具体模型组成的模型池，并将其分为四个能力层级（表~\ref{tab:appendix-model-pool}）。目标层级标签相对于该模型池作存在性定义：在 \S\ref{sec:dataset-construction} 的降级--级联过程中，只要层级 $t$ 的模型池中任一个模型能够正确处理当前步骤（依据完整轨迹仍以相同步数通过来推断），便接受该记录属于层级 $t$。公开静态记录只存储得到的层级标签，而不存储具体模型 ID，因此当底层模型池更新时，基于这些数据训练的路由器仍保持可移植性。

\begin{table}[ht]
\centering
\small
\caption{静态模型池（11 个模型，四个层级）。同一逻辑模型接受的别名在括号中显示。动态赛道代表模型的定价单列于表~\ref{tab:appendix-pricing}。}
\label{tab:appendix-model-pool}
\begin{adjustbox}{max width=\linewidth}
\begin{tabular}{lp{0.65\linewidth}}
\toprule
层级 & 模型 \\
\midrule
\texttt{high}     & \texttt{anthropic/claude-opus-4.6}（别名：
                     \texttt{anthropic/claude-opus-4-6}）；
                     \texttt{openai/gpt-5.4}（别名：
                     \texttt{openai/gpt-5.4-2026-03-05}） \\
\texttt{mid\_high} & \texttt{anthropic/claude-haiku-4-5}；
                     \texttt{google/gemini-3-flash-preview}（别名：
                     \texttt{gemini/gemini-3-flash-preview}）；
                     \texttt{qwen/qwen3.5-397b-a17b} \\
\texttt{mid}      & \texttt{minimax/minimax-m2.5}；
                     \texttt{qwen/qwen3.5-27b}；
                     \texttt{qwen/qwen3-coder} \\
\texttt{low}      & \texttt{deepseek/deepseek-v3.2}；
                     \texttt{z-ai/glm-4.5-air}；
                     \texttt{qwen/qwen3.5-9b} \\
\bottomrule
\end{tabular}
\end{adjustbox}
\end{table}

\subsection{静态层级计费费率}
\label{app:static-tier-rates}

\begin{table}[ht]
\centering
\small
\begin{tabular}{lrrrr}
\toprule
 & 输入 & 缓存读取 & 缓存写入 & 输出 \\
层级 & （\$/1M） & （\$/1M） & （\$/1M） & （\$/1M） \\
\midrule
low & 0.26 & 0.13 & 0.26 & 0.5 \\
mid & 0.30 & 0.059 & 0.30 & 2.0 \\
mid\_high & 0.50 & 0.05 & 0.083 & 5.0 \\
high & 5.0 & 0.50 & 6.25 & 25.0 \\
\bottomrule
\end{tabular}
\caption{静态评分协议的逐层级计费常数（美元/百万 token）。这些费率并非任一单独模型的价格，而是从各层级多个模型的价格区间中得出的代表值。\texttt{high} 层级比其余三个层级昂贵 10--50$\times$，后三者之间的价格仅相差 2--5$\times$。动态模型池价格见表~\ref{tab:appendix-pricing}。这种不对称驱动附录~\ref{app:cost-ablation}中的消融。}
\label{tab:pricing}
\end{table}

\subsection{动态模型池的代表模型与定价}
\label{app:pricing}

\begin{table}[ht]
\centering
\small
\caption{TwinRouterBench 动态模型池的具体代表模型及实时 OpenRouter 单价。这些模型价格不同于表~\ref{tab:pricing}中的静态层级计费常数。价格以美元/百万 token 表示；当提供商未公布单独的缓存写入费率时，缓存写入价格回退为输入价格。这些是动态赛道评分时替入 $c_i(t)$ 的具体 $p$ 值；发布的评分器按包含固定未解决惩罚成本（\S\ref{sec:dynamic-scoring}）的排行榜账单排序。选择 \texttt{mid} 层的代表模型 \texttt{minimax-m2.7}，是因为它在 SWE-bench 排行榜上同等价位模型中排名最高，并与 \texttt{minimax-m2.5}（静态标签构建中使用的中层代表）处于相同层级；表~\ref{tab:appendix-model-pool}的静态池冻结在构建层级标签时所用的版本。}
\label{tab:appendix-pricing}
\begin{adjustbox}{max width=\linewidth}
\begin{tabular}{llllrrrr}
\toprule
层级 & 代表模型 & 上下文 & 最大输出 & 输入 & 输出 & 缓存读取 & 缓存写入 \\
\midrule
high     & anthropic/claude-opus-4.6       & 1M      & 128K   & 5.00  & 25.00 & 0.50   & 6.25  \\
mid\_high & google/gemini-3-flash-preview  & 1.05M   & 65.5K  & 0.50  &  3.00 & 0.05   & 0.0833 \\
mid      & minimax/minimax-m2.7             & 196.6K  & 196.6K & 0.30  &  1.20 & 0.059  & 0.30  \\
low      & deepseek/deepseek-v3.2           & 163.8K  & 163.8K & 0.252 &  0.378 & 0.0252 & 0.252 \\
\bottomrule
\end{tabular}
\end{adjustbox}
\end{table}

\subsection{规则式 UncommonRoute 配置}
\label{app:ur-config}

规则式 UncommonRoute 在双信号集成中使用 \texttt{MetadataSignal} 与 \texttt{EmbeddingSignal}，采用上游默认值 \path{risk_tolerance}$=0.5$、\path{ensemble_weights}$=(0.55,0.45)$ 和 \path{low_escalation_threshold}$=0.55$。由于该规则式设置中不存在基准特定的种子集、分类器和 Platt scaler，嵌入信号会弃权，实践中只有元数据信号产生贡献。

\subsection{步级成本案例研究：\texttt{sympy\_\_sympy-12096}}
\label{app:sympy}

表~\ref{tab:appendix-sympy}展示一条 13 步 SWE-bench 轨迹。方案 A 是带提示词缓存的全 \texttt{high} 执行；方案 B 是经验证的步级路由（层级分配见“层级”列）。步级路由在解决该 issue 的同时，相对全 \texttt{high} 实现了 39.6\% 的成本降低。最后一步的输出 token 为估计值~($*$)。

\begin{table}[ht]
\centering
\scriptsize
\caption{\texttt{sympy\_\_sympy-12096} 的步级成本案例研究。方案 A = 带提示词缓存的全 \texttt{high}。方案 B = 经验证的步级路由。token 数量以千计；成本单位为美元。}
\label{tab:appendix-sympy}
\begin{adjustbox}{max width=\linewidth}
\begin{tabular}{r l r r r r r r}
\toprule
步骤 & 层级 & 方案 A 输入 & 方案 A 输出 & 方案 A 成本 & 方案 B 输入 & 方案 B 输出 & 方案 B 成本 \\
\midrule
 1 & mid      & 1,321 & 112 & \$0.0111 & 1,284 & 108 & \$0.0005 \\
 2 & mid      & 1,744 &  71 & \$0.0051 & 1,699 &  68 & \$0.0003 \\
 3 & mid\_high & 1,846 &  69 & \$0.0032 & 1,647 &  66 & \$0.0003 \\
 4 & mid\_high & 2,502 &  67 & \$0.0067 & 2,343 &  65 & \$0.0003 \\
 5 & low      & 3,287 &  89 & \$0.0084 & 3,729 &  87 & \$0.0010 \\
 6 & mid\_high & 4,103 & 256 & \$0.0131 & 3,900 & 254 & \$0.0010 \\
 7 & low      & 4,512 &  55 & \$0.0060 & 5,215 &  55 & \$0.0009 \\
 8 & mid\_high & 4,612 & 168 & \$0.0071 & 4,403 & 168 & \$0.0007 \\
 9 & high     & 4,921 & 241 & \$0.0103 & 4,921 & 241 & \$0.0368 \\
10 & high     & 5,149 &  85 & \$0.0060 & 5,149 &  85 & \$0.0060 \\
11 & high     & 5,591 &  63 & \$0.0069 & 5,591 &  63 & \$0.0069 \\
12 & low      & 5,653 &  56 & \$0.0046 & 6,789 &  59 & \$0.0018 \\
13 & low      & 5,864 & 111$^{*}$ & \$0.0069 & 7,077 & 109$^{*}$ & \$0.0010 \\
\midrule
总计 & -- & -- & -- & \textbf{\$0.0953} & -- & -- & \textbf{\$0.0576} \\
\bottomrule
\end{tabular}
\end{adjustbox}
\end{table}

\subsection{mtRAG/QMSum 裁判加固细节}
\label{app:judge}

表~\ref{tab:appendix-judge}记录了从旧版裁判准则到最终数据集所用加固版本的变更。每一行都指明旧版裁判中一种具体的伪通过失败模式，以及构建过程中采用的对应加固规则。

\begin{table}[ht]
\centering
\small
\caption{mtRAG/QMSum 裁判加固。这些改动将薄弱的参考相似度判断转化为保守的任务成功谓词，使其符合当前轮/查询要求和 RadBench 风格的忠实性标准~\citep{kuo2025radbench}。}
\label{tab:appendix-judge}
\begin{tabularx}{\linewidth}{p{0.22\linewidth}p{0.35\linewidth}X}
\toprule
旧版裁判的风险 & 加固规则 & 目的 \\
\midrule
缺乏首要准则 &
  先判定答案是否完整且正确地解决当前轮/查询 &
  防止答案与参考的表面相似替代任务成功 \\
将参考当作答案字符串 &
  将参考视为覆盖提示；证据/转录文本具有权威性 &
  允许有效改述，并拒绝盲从参考的错误 \\
无支持事实 &
  对没有支持的关键数字、日期、金额、比例、投票或主张直接硬失败 &
  抑制看似合理的幻觉答案 \\
证据冲突 &
  当候选答案与检索段落或会议转录冲突时直接硬失败 &
  强制忠实性 \\
不可回答/无上下文样例 &
  mtRAG 预过滤：必须可回答、具有上下文段落且有人类正向相关性反馈 &
  在路由搜索前移除结构上不可靠的样例 \\
所有层级均失败 &
  标记为丢弃；不生成样例 JSON 或监督记录 &
  防止伪通过成为已验证标签 \\
\bottomrule
\end{tabularx}
\end{table}

\subsection{前沿 LLM-as-Router 诊断汇总}
\label{app:opus}

表~\ref{tab:opus-high}展示 Opus~4.6 对经验证的 \texttt{high} SWE-bench 步骤的预测；表~\ref{tab:appendix-opus-bias}提供额外诊断。这些数值支持如下观点：强模型未经执行的路由判断，不能可靠替代经执行验证的步级监督（\S\ref{sec:opus}）。

\begin{table}[ht]
\centering
\small
\begin{tabular}{lrrrrr}
\toprule
验证层级 & $\tilde{t}=\texttt{low}$ & $\tilde{t}=\texttt{mid}$ & $\tilde{t}=\texttt{mid\_high}$ & $\tilde{t}=\texttt{high}$ & 总计 \\
\midrule
$\hat{t}=\texttt{high}$ & 34 & 43 & 63 & \textbf{7} & 147 \\
\bottomrule
\end{tabular}
\caption{Opus 4.6 对经验证 \texttt{high} SWE-bench 步骤的预测。147 个步骤中有 140 个被路由至低于 \texttt{high} 的层级。}
\label{tab:opus-high}
\end{table}

\begin{table}[ht]
\centering
\small
\caption{SWE-bench 上前沿 LLM-as-router（Opus~4.6）的诊断汇总。核心发现是对 \texttt{high} 层级步骤的严重低估：仅 5\% 的经验证 \texttt{high} 步骤被预测为 \texttt{high}，导致该路由下每一条 SWE-bench 轨迹均失败。}
\label{tab:appendix-opus-bias}
\begin{tabularx}{\linewidth}{lrr X}
\toprule
诊断项 & 数值 & 范围 & 解释 \\
\midrule
SWE-bench 有效记录 & 300 & 共 336 条记录，36 条格式错误 & 用于混淆矩阵分析的有效子集 \\
验证为 high、预测为 high & 7/147（4.8\%） & SWE-bench 有效记录 & 对 high 层级步骤的严重低估 \\
验证为 high、被低估 & 140/147（95.2\%） & SWE-bench 有效记录 & 大多数 high 步骤被路由到 low/high 以下的层级 \\
SWE-bench 轨迹通过 & 0/40 & 旧版 Opus 路由运行 & 每条轨迹都有 $\ge$1 个路由不足的步骤 \\
中间三分之一位置的低估 & 73\% & 步骤位置诊断 & 核心编辑/测试运行步骤最易被低估 \\
首个三分之一位置的低估 & 46\% & 步骤位置诊断 & 程度较低但仍有显著的路由不足 \\
末个三分之一位置的低估 & 53\% & 步骤位置诊断 & 恢复/收尾步骤仍然困难 \\
\bottomrule
\end{tabularx}
\end{table}

\subsection{标签验证审计轨迹}
\label{app:validation}

表~\ref{tab:appendix-validation}记录数据集构建期间采用的验证机制。它们共同确保，除非强模型基线通过、层级池中的模型确实在混合模型轨迹下通过、裁判未将答案标为伪通过，且最终人工审计未发现遗留问题，否则没有标签会进入发布的 JSONL。

\begin{table}[ht]
\centering
\small
\caption{标签验证审计轨迹。每个组成部分针对构建流水线中的一种特定失败模式。}
\label{tab:appendix-validation}
\begin{tabularx}{\linewidth}{p{0.22\linewidth}p{0.22\linewidth}p{0.28\linewidth}X}
\toprule
验证组件 & 范围 & 处理的失败模式 & 发布操作 \\
\midrule
成功种子轨迹 &
  全部发布实例 &
  将任务失败与路由失败混淆 &
  只有通过的种子轨迹才进入降级搜索 \\
顺序锁定降级搜索 &
  多步轨迹 &
  漏掉最后一个通过层级；在最强轨迹上叠加假设性输出 &
  在因果的混合模型前缀下，从 Opus 建议之下推一层，直至 harness 失败，再锁定最后一个通过层级 \\
层级级联（每层 3 个模型） &
  全部非 \texttt{low} 层级标签 &
  同一层级中单模型的假阴性 &
  只有一次实际运行通过后才锁定更低成本层级 \\
严格的 mtRAG/QMSum 裁判 &
  开放式 RAG 与摘要 &
  薄弱判定产生的伪通过 &
  使用硬失败准则；丢弃全部层级失败的样例 \\
最终人工审计 &
  BFCL、SWE-bench 和 PinchBench 各自约 $\sim$10\% 的随机步骤样本 &
  上游搜索没有收紧到条件最优的、仍可降级的层级标签 &
  人工审阅在路由器可见前缀下决定发布层级能否再降一级；一个 SWE-bench 步骤被标记为可继续降级，并相应下调其层级 \\
\bottomrule
\end{tabularx}
\end{table}

\subsection{人工审计协议与汇总}
\label{app:manual-audit}

人工审计是静态赛道的最后一道质量关：它在由 Opus 初始化的降级搜索、层级池级联检查、开放式裁判加固，以及 \S\ref{sec:dataset-construction} 的前缀重建都完成\emph{之后}执行。审计有意采用人工判断而非第二次 harness 执行。其目的不是重新计算任何通过谓词，而是在部署路由器也会看到的同一份路由器可见上下文中，交叉检查已发布层级标签是否仍然站得住脚。

\paragraph{抽样。}
我们遵循发布制品中包含的 GT 层级最优性审查协议。审计单元是一条被路由的步骤，使单步和多步样例在相同粒度上处理。在含有多步轨迹的三类工作负载（BFCL、SWE-bench、PinchBench）中，我们从该工作负载的路由步骤独立随机抽取约 10\%。mtRAG 和 QMSum 是缺乏多轮累积复杂度的单轮任务；其标签已经受到加固裁判的约束，伪通过模式已通过迭代校准消除（\S\ref{sec:dataset-construction}），因此不在此审计中重新抽样。

\paragraph{审阅流程。}
对每个抽样步骤，将路由器可见的 \texttt{messages} 前缀固定为发布版本。可附带一个可选的大模型预标签供参考，但最终决定始终由人工审阅者作出；审阅者判断将已验证层级再降一级后，是否仍会产生正确处理当前步骤的响应。步骤随后标记为\emph{紧致}、\emph{不确定}或\emph{可继续降级}。\emph{紧致}判定表示发布层级已处于条件最优，无法再降一级而不破坏任务解决；\emph{可继续降级}判定表示审阅者认为发布层级仍然过高，标签可以再收紧一级。审计总共检查 64 个步骤；一个 SWE-bench 步骤被标记为可继续降级，且其验证层级在发布前下调一级，另 63 个步骤被判定为紧致。这一内部审计旨在捕获上游搜索未收紧至条件最优的残余标签，并澄清标签语义；它并不被表述为外部标注研究，也不是完整语料库的置信区间。

\begin{table}[ht]
\centering
\small
\caption{人工审计汇总。当审阅者认为已验证层级还能再降低一级而不破坏任务解决时，抽样步骤被计为\emph{可降级}；\emph{紧致比例}是其补集，即层级标签已处于条件最优的比例。抽样在含多步轨迹的三个工作负载上按步骤粒度进行；mtRAG 和 QMSum 是已由加固裁判保护的单轮任务，此处不重新抽样。该审计是有针对性的内部质量检查。}
\label{tab:manual-audit}
\begin{tabular}{lrrrrr}
\toprule
切片 & 总步骤数 & 抽样步骤数 & 抽样比例 \% & 可降级 & 紧致比例 \% \\
\midrule
BFCL & 248 & 25 & 10.1 & 0 & 100.0 \\
SWE-bench & 336 & 34 & 10.1 & 1 & 97.1 \\
PinchBench & 48 & 5 & 10.4 & 0 & 100.0 \\
\midrule
总计 & 632 & 64 & 10.1 & 1 & 98.4 \\
\bottomrule
\end{tabular}
\end{table}

\subsection{静态赛道分解}
\label{app:static-breakdown}

\paragraph{按工作负载划分的成本节省。}

表~\ref{tab:per-benchmark}显示 SWE-bench 是最困难的静态切片。它在按记录加权的 \textsc{CostSave} 分数中占 34.64\%；由于一个路由不足的步骤就会使整条轨迹失败并触发考虑失败的账单，每个路由器在该切片上的成本节省均为负。对于其余四类工作负载，SR-KNN 和 ClawRouter（规则）平均均超过 85\% 的成本节省；全局排名取决于路由器是否保持代码修复轨迹。

\begin{table}[ht]
\centering
\scriptsize
\begin{tabular}{lrrrrrr}
\toprule
路由器 & BFCL & mtRAG & QMSum & Pinch & SWE-bench & 总体 \\
\midrule
SR-KNN       & \textbf{90.49} & \textbf{92.35} & 90.24 & 35.36 & \textbf{-1.64} & \textbf{56.18} \\
ClawRouter（规则） & 89.87 & 83.45 & \textbf{92.12} & \textbf{55.20} & -6.55 & 53.82 \\
UncommonRoute（规则） & 47.14 & 91.93 & 90.24 & 39.80 & -86.08 & 15.99 \\
\midrule
记录权重  & 25.57 & 19.90 & 14.95 & 4.95 & 34.64 & 100.00 \\
\bottomrule
\end{tabular}
\caption{静态评分协议下按工作负载分解的 \textsc{CostSave}。SWE-bench 切片是主要压力测试：失败轨迹会触发 \S\ref{sec:static-scoring} 中考虑失败的分子。}
\label{tab:per-benchmark}
\end{table}

\paragraph{路由器失败模式。}
\label{sec:error-patterns}

三个基线以不同方式失败（表~\ref{tab:error-patterns}）。在 800 条验证层级低于 \texttt{high} 的记录中，SR-KNN 大多数情况下精确匹配；ClawRouter（规则）几乎总是比标签高一层；UncommonRoute（规则）的精确匹配比 ClawRouter（规则）更多，但会 217 次跳至 \texttt{high}。

\begin{table}[ht]
\centering
\small
\begin{tabular}{lrrr}
\toprule
$\hat{t}<\texttt{high}$ 时的预测模式 & SR-KNN & ClawRouter（规则） & UncommonRoute（规则） \\
\midrule
精确层级                         & \textbf{627} & 107 & 488 \\
高一层                         & 10 & \textbf{649} & 47 \\
跳至 \texttt{high}               & 117 & 21 & \textbf{217} \\
失败：预测低于 $\hat{t}$     & 46 & 23 & \textbf{48} \\
\bottomrule
\end{tabular}
\caption{800 条非 \texttt{high} 记录上的错误模式。}
\label{tab:error-patterns}
\end{table}

SWE-bench 说明了为何记录级行为并不充分（表~\ref{tab:swe-high}）。UncommonRoute（规则）识别出许多经验证的 \texttt{high} 步骤，但在一条包含 8--13 次调用的轨迹中即使漏掉一个步骤，也会使该实例失败。

\begin{table}[ht]
\centering
\small
\begin{tabular}{lrrr}
\toprule
路由器 & 在 $\hat{t}=\texttt{high}$ 时预测为 \texttt{high} & 命中率 & 失败轨迹 \\
\midrule
SR-KNN       & 137/168 & \textbf{81.5\%} & \textbf{10/40} \\
ClawRouter（规则） & 7/168   & 4.2\%  & 38/40 \\
UncommonRoute（规则） & 105/168 & 62.5\% & 39/40 \\
\bottomrule
\end{tabular}
\caption{SWE-bench 的高层级决策。一个路由不足的关键步骤就会使整条轨迹失败。}
\label{tab:swe-high}
\end{table}

\subsection{成本节省计费消融}
\label{app:cost-ablation}

本节给出 \S\ref{sec:main-results}引用的完整推导和逐变体数值。

\paragraph{定价不对称性驱动消融。}
表~\ref{tab:save-asymmetry}针对三个代表性步骤配置，报告各目标层级相对于始终 \texttt{high} 基线的逐步节省。对于三种步骤规模，\texttt{save\_mid / save\_low} 比值都落在 $1$ 的 3\% 以内，这意味着路由器从 $\hat{t}=\texttt{low}$ 错路由到 $\tilde{t}=\texttt{mid}$ 在绝对金额上几乎不损失成本。

\begin{table}[ht]
\centering
\small
\caption{三个代表性步骤配置上，不同目标层级相对于始终 \texttt{high} 的节省。从 \texttt{low} 错路由到 \texttt{mid} 几乎没有成本。}
\label{tab:save-asymmetry}
\begin{tabular}{lrrrr}
\toprule
步骤配置 & 节省（high$-$low） & 节省（high$-$mid） & 节省（high$-$mid\_high） & save\_mid / save\_low \\
\midrule
4k 冷步骤     & 3.621¢ & 3.530¢ & 3.467¢ & 0.975 \\
20k 热步骤    & 1.965¢ & 2.032¢ & 1.900¢ & 1.034 \\
100k 冷步骤   & 61.13¢ & 60.65¢ & 62.67¢ & 0.992 \\
\bottomrule
\end{tabular}
\end{table}

\paragraph{五种评分变体。}
我们在同一组预测上比较五种 $N/D$ 计费方案：

\begin{itemize}
\item \textbf{A}（旧版记录级分母）：对满足 $\hat{t}_i \ne \texttt{high}$ 的通过记录，$D = \sum_i (\mathrm{cost}(3;i) - \mathrm{cost}(\hat{t}_i;i))$；$N$ 忽略失败。
\item \textbf{B}（旧版全记录分母）：$D$ 与 A 相同，但覆盖全部记录；失败轨迹在轨迹层面减去 $\sum_i \mathrm{cost}(\tilde{t}_i;i)$。
\item \textbf{C}（中间的 high 记录排除）：从 $D$ 中排除 $\hat{t}_i=\texttt{high}$ 的记录。
\item \textbf{D}（$D = \sum_i \mathrm{cost}(3;i)$，在步骤层面惩罚失败，没有轨迹级惩罚）。
\item \textbf{E}（最终的考虑失败公式）：$D$ 与 D 相同，但失败轨迹在轨迹层面收取 $-\sum_i \mathrm{cost}(\tilde{t}_i;i)$。
\end{itemize}

\begin{table}[ht]
\centering
\small
\caption{\textsc{CostSave} 计费消融。变体 A--D 均将 ClawRouter（规则）排在 SR-KNN 之上；只有变体 E（最终的考虑失败公式）与 \textsc{RowPass}、\textsc{RowExact} 和 \textsc{TrajPass} 一致。}
\label{tab:ablation-scoring}
\begin{tabular}{lrrr}
\toprule
变体 & SR-KNN \textsc{cost}\% & ClawRouter（规则）\textsc{cost}\% & 排名 \\
\midrule
A（旧版记录级分母）                  & 84.99 & \textbf{89.71} & Claw 胜 \\
B（旧版全记录分母）                    & 84.45 & \textbf{84.59} & Claw 胜 \\
C（中间 high 记录排除）               & 79.11 & \textbf{91.31} & Claw 胜 \\
D（$D=\sum\mathrm{cost}(3)$，步骤级失败）    & 61.81 & \textbf{67.67} & Claw 胜 \\
\textbf{E}（最终考虑失败公式）          & \textbf{56.18} & 53.82 & SR 胜 \\
\bottomrule
\end{tabular}
\end{table}

\paragraph{为什么变体 E 恢复一致的排名。}
变体 E 将一条失败轨迹解释为“路由器烧掉了自己的（低价）预算，\emph{并且}用户支付完整 \texttt{high} 账单来重新计算”。形式上，失败轨迹上用户侧账单为 $\sum_i \mathrm{cost}(\tilde{t}_i;i) + \sum_i \mathrm{cost}(3;i)$，故相对于始终 \texttt{high} 的节省等于 $-\sum_i \mathrm{cost}(\tilde{t}_i;i)$，这正是 $N$ 所编码的值。重试项 $\sum_i \mathrm{cost}(3;i)$ 被 $D$ 吸收而非重复计数，使 \textsc{CostSave} $\in [-\infty, 100]$，其中 $0$ 代表始终 \texttt{high} 的路由。在完整基准上，ClawRouter（规则）有 649 个高一层错误（在绝对金额上近乎免费），这些错误不再能弥补其 38 条失败的 SWE-bench 轨迹，因此其 \textsc{CostSave} 与 \textsc{RowPass}/\textsc{RowExact}/\textsc{TrajPass} 数值一致。

\subsection{Opus 4.6 混淆矩阵}
\label{app:opus-confusion}

供参考，在先前的三指标计费协议下记录的 Opus 4.6 路由运行，整体分数为 84.05，其中记录通过 $=$85.77、记录精确 $=$82.88、成本节省 $=$83.48。它不能与表~\ref{tab:main-results}的静态分数直接比较，因为先前的计费并未使用 \S\ref{sec:static-scoring} 中考虑失败的轨迹惩罚，也未报告轨迹通过指标；因此我们将其视为历史数据点而非用于排名的基线。

表~\ref{tab:opus-confusion}展示了 \S\ref{sec:opus}中汇总的完整混淆矩阵。

\begin{table}[ht]
\centering
\small
\caption{300 条有效 SWE-bench 记录（排除 36 条格式错误输出）上，Opus 4.6 的路由预测与验证层级。对角单元格加粗。即使验证层级为 \texttt{high}，Opus 也很少选择 \texttt{high}。}
\label{tab:opus-confusion}
\begin{tabular}{lrrrrr}
\toprule
 & $\tilde{t}=\texttt{low}$ & $\tilde{t}=\texttt{mid}$ & $\tilde{t}=\texttt{mid\_high}$ & $\tilde{t}=\texttt{high}$ & 行总计 \\
\midrule
$\hat{t}=\texttt{low}$       & \textbf{52（61\%）} & 15 & 17 & 1 & 85 \\
$\hat{t}=\texttt{mid}$       & 13 & \textbf{7（23\%）} & 10 & 0 & 30 \\
$\hat{t}=\texttt{mid\_high}$ & 8  & 13 & \textbf{16（42\%）} & 1 & 38 \\
$\hat{t}=\texttt{high}$      & 34 & 43 & 63 & \textbf{7（5\%）} & 147 \\
\bottomrule
\end{tabular}
\end{table}

\subsection{mtRAG 记录示例}
\label{app:mtrag-example}

除 \S\ref{sec:benchmark-overview} 中的 SWE-bench 示例外，第二种记录格式说明 \texttt{low} 层级标签如何产生：经过若干次检索轮次后，一个简短的后续查询可以由模型池中最低成本层级回答。

\begin{center}
\begin{minipage}{0.95\linewidth}
\scriptsize
\begin{verbatim}
{
  "id": "mtrag_..._turn_3_step_1",
  "benchmark": "mtrag",
  "step_index": 1, "total_steps": 1,
  "messages": [
    {"role": "system", "content": "Answer based on the retrieved passages..."},
    {"role": "user",   "content": "[Passage] Major.minor update... [Passage] ..."},
    {"role": "assistant", "content": "..."},
    ...multi-turn history...
    {"role": "user",   "content": "Major.minor update."}
  ],
  "target_tier": "low", "target_tier_id": 0
}
\end{verbatim}
\end{minipage}
\end{center}
